\chapter{Introduction}
\label{chap:introduction}

\section{Motivation}

Post-COVID monetary policy brought an older question back into sharp focus: monetary transmission is not only about the average response of inflation or output, but also about the channels through which financial conditions, bank intermediation, housing credit, and compensation-linked variables adjust over time. The Euro Area entered the pandemic after a long period of low rates, unconventional policy, and balance-sheet expansion. It then experienced extraordinary accommodation, followed by a rapid tightening cycle as inflation accelerated. In that environment, the same ECB announcement can move market prices inside the high-frequency window, shift lending margins over subsequent months, affect house-purchase credit growth, and leave a weaker or less consistent trace in compensation-linked proxies.

This thesis studies that heterogeneous transmission. Its central argument is not that monetary policy mainly inflated assets or that broad real-income transmission was absent. Those claims exceed what aggregate local projections can identify. The thesis instead asks whether post-COVID ECB easing shocks transmitted most clearly and persistently through housing-credit and selected financial-intermediation variables, while evidence for broad real-income transmission was weaker and less consistent.

The empirical contribution is therefore mechanism-focused. The required interpretation chain is:

\begin{figure}[htbp]
  \centering
  \fbox{%
    \begin{minipage}{0.92\textwidth}
      \centering
      ECB policy-news shock
      \(\rightarrow\) lending conditions and intermediation variables
      \(\rightarrow\) house-purchase credit growth
      \(\rightarrow\) selected financial and asset-market responses
    \end{minipage}%
  }
  \caption{Reduced-form transmission chain guiding the empirical interpretation. The diagram is a mechanism guide, not a structurally identified mediation model.}
  \label{fig:conceptual-intermediation-chain}
\end{figure}

\noindent The model is intentionally not redesigned around a new structural architecture. It remains a high-frequency identified, dynamic impulse-response framework. The revision is interpretive: housing-credit transmission and intermediation variables become the center of the thesis, while equity prices and broad financial indices become supporting evidence that must be read with information-effect caution.

\section{Research Question}

The thesis is organized around the following research question:

\begin{quote}
Did post-COVID ECB monetary easing transmit more strongly through housing-credit and financial-intermediation channels than through compensation-linked real-economy variables?
\end{quote}

The question is comparative, dynamic, and reduced-form. It concerns relative response strength and persistence across observable channels. It does not ask whether monetary policy caused household-level welfare changes, generated inequality, or shifted resources across population groups. Those questions require household, regional, or bank-level microdata that are outside the scope of the empirical design.

\section{Hypotheses}

Three hypotheses discipline the empirical interpretation.

\begin{hypothesis}[Housing-credit persistence]
\label{hyp:housing-credit-persistence}
Expansionary ECB policy-news shocks generated clear and persistent responses in housing-credit variables, especially house-purchase lending growth.
\end{hypothesis}

\begin{hypothesis}[Compensation-linked attenuation]
\label{hyp:compensation-linked-attenuation}
Compensation-linked real-economy variables responded less consistently and less persistently than the housing-credit block.
\end{hypothesis}

\begin{hypothesis}[Differential financial-intermediation transmission]
\label{hyp:differential-financial-intermediation}
Expansionary monetary-policy shocks generated stronger and more persistent responses in housing-credit and financial-intermediation variables than in productive-credit and compensation-linked variables, indicating heterogeneous monetary transmission across sectors.
\end{hypothesis}

These hypotheses keep the thesis close to the evidence. \Cref{hyp:housing-credit-persistence} gives house-purchase lending growth its central role. \Cref{hyp:compensation-linked-attenuation} frames wages and labor-income proxies as a comparative block rather than as a welfare test. \Cref{hyp:differential-financial-intermediation} replaces overbroad distributional language with a testable macro-financial transmission claim.

\section{Empirical Object}

The thesis estimates monthly local-projection responses to high-frequency ECB target-factor surprises aggregated to the month and signed so positive values denote easing. The canonical shock is \nolinkurl{target_factor_monthly_easing}. Responses are estimated at h0, h1, h3, h6, h12, and h24 and normalized to a one-standard-deviation surprise.

The empirical object is a horizon-specific response profile to measured ECB policy news. The comparison spans house-purchase lending growth, pure-new mortgage flows, mortgage and NFC lending spreads, broad credit stocks, ECB liquidity-stock context, equity-market responses, labor expectations, and negotiated wage-pressure proxies. Quarterly house prices, compensation per employee, and Bank Lending Survey indicators are not interpolated into monthly responses; they remain context or robustness material when their observed frequency allows.

This object supports claims about differential transmission strength, response persistence, timing, and robustness. It does not identify household-level incidence, welfare effects, inequality, redistribution, exact structural mediation, or bank-specific treatment effects. That boundary is central to the thesis because the final inference is strongest when the estimand remains tied to observable repository outputs.

\section{DAX Sign and Information Effects}

The real DAX response is a central interpretation boundary. A positive easing surprise is followed by a negative real DAX response at medium and longer horizons in the baseline results. That pattern contradicts a simple story in which easing mechanically inflates asset prices. The equity evidence must therefore be read through the information-effects literature: ECB announcements can reveal central-bank information about future macroeconomic conditions as well as policy accommodation.

In this interpretation, a negative equity response may indicate that high-frequency easing surprises partly contain adverse macroeconomic information revealed during ECB communication. Markets react to both policy and information. Easing can therefore coincide with deteriorating growth expectations, weaker expected earnings, or changes in risk premia. This does not make the DAX result irrelevant. It makes it a cautionary financial-market result rather than evidence of broad asset-price appreciation.

\section{Contribution}

The thesis contributes a disciplined macro-financial transmission study in five ways.

\begin{enumerate}
  \item It integrates high-frequency ECB policy-news surprises into a monthly local-projection framework and traces dynamic responses across macro-financial channels.

  \item It makes house-purchase lending growth the central empirical variable. The strongest result is not a broad asset-market claim, but a housing-credit response with clear relative persistence.

  \item It explicitly models financial intermediation through lending spreads, lending rates, credit aggregates, balance-sheet context, and timing evidence. These variables provide the bridge between policy news and housing-credit expansion.

  \item It interprets equity-market responses cautiously. The negative real DAX response is treated as mixed financial evidence consistent with possible information effects, not as a contradiction to be hidden.

  \item It states the claim boundary directly. The thesis provides reduced-form evidence of heterogeneous monetary transmission; it is not a structural welfare analysis, inequality decomposition, or household distribution study.
\end{enumerate}

\section{Structure of the Thesis}

\Cref{chap:literature} reviews monetary transmission theory, high-frequency identification and information effects, housing-credit transmission, and restrained financialization arguments. \Cref{chap:theory} develops the mechanism framework. \Cref{chap:methodology} defines the shock, estimator, uncertainty layer, proxy governance, and timing logic. \Cref{chap:results} reports identification validity, housing-credit transmission, intermediation responses, equity and asset-market evidence, compensation-linked outcomes, and regime comparisons. \Cref{chap:robustness} tests the durability of the hierarchy. \Cref{chap:discussion} interprets the surviving evidence within reduced-form boundaries. \Cref{chap:conclusion} concludes.
